\SourceLeaf{216}%
\bigskip\Needspace{10\baselineskip}
\pdfbookmark[1]{Fortgesetzte Nachrichten. Ceres Ferdinandea. Mai 1802}{ceres180205}
\begin{center}
{\small FORTGESETZTE NACHRICHTEN ÜBER DEN NEUEN HAUPT-PLANETEN.\par}
CERES FERDINANDEA.\par
\medskip\hrule\medskip
Monatliche Correspondenz. Band V. S.\ 462—481. Mai 1802.\par
\medskip\hrule
\end{center}\medskip
—\quad —\quad —\quad Für die Zeit bis zum 3. März haben wir im vorigen Heft auf S.\ 389 [S.\ 209 dieses Bandes] die Übereinstimmung unserer Beobachtungen mit den Elementen von \textsc{Gauß} angegeben. Seitdem hat Dr.\ \textsc{Gauß} den Vergleich fortgesetzt, wie hier zu sehen ist.
\begin{center}
{\small\setlength{\tabcolsep}{5pt}\renewcommand{\arraystretch}{0.92}
\begin{tabular}{r r|r|r|r|r}
\multicolumn{2}{c|}{\shortstack{Seeberg\\1802}}&\shortstack{Berechnete\\R.A. der $\CeresSymbol$}&\shortstack{Unter-\\schied}&\shortstack{Berechnete\\Dekl. der $\CeresSymbol$}&\shortstack{Unter-\\schied}\\\hline
März&6&$185^0\ 49'\ 23''3$&$+4''6$&$16^0\quad4'\ 12''2$&$+23''0$\\
—&7&$185\ 38\ 22.0$&$+6.1$&$16\ 10\ 45.6$&$+29.7$\\
—&10&$185\quad3\ 57.2$&$+8.3$&$16\ 29\ 48.4$&$+29.5$\\
—&11&$184\ 52\quad5.2$&$+9.4$&$16\ 35\ 55.2$&$\ .\ .\ .\ .\ .$\\
—&15&$184\quad3\quad2.9$&$+10.9$&$16\ 58\ 58.2$&$+27.5$\\
—&16&$183\ 50\ 28.4$&$+9.4$&$17\quad4\ 20.6$&$+21.7$\\
—&17&$183\ 37\ 48.5$&$+13.7$&$17\quad9\ 32.6$&$+26.2$\\
—&18&$183\ 25\quad3.2$&$+12.5$&$17\ 14\ 34.4$&$+29.1$\\
—&19&$183\ 12\ 14.4$&$+15.8$&$17\ 19\ 25.1$&$+34.3$
\end{tabular}}
\end{center}

Dr.\ \textsc{Gauß} verglich auch zwei \textit{Greenwicher} Beobachtungen mit seinen Elementen, und die Rektaszensionen stimmen aufs Genaueste mit unseren überein, wenn man, wie oben, als unzweifelhaft voraussetzt, dass bei der Beobachtung vom 19. Februar die Minute versehentlich falsch notiert wurde.
\begin{center}
{\small\setlength{\tabcolsep}{2.8pt}\renewcommand{\arraystretch}{0.92}
\begin{tabular}{l r|r|r|r|r|r|r}
\multicolumn{2}{c|}{}&&\multicolumn{2}{c|}{Unterschied bei}&&\multicolumn{2}{c}{Unterschied bei}\\
\multicolumn{2}{c|}{\shortstack{Greenwich\\1802}}&\shortstack{Berechnete\\R.A. der $\CeresSymbol$}&\shortstack{Dr.\ \textsc{Mas-}\\\textsc{kelyne}}&\textsc{v. Zach}&\shortstack{Berechnete\\Deklination\\der $\CeresSymbol$}&\shortstack{Dr.\ \textsc{Mas-}\\\textsc{kelyne}}&\textsc{v. Zach}\\\hline
Februar&19&$187^0\ 58'\ 13''3$&$-4''3$&$-4''1$&$14^0\ 20'\ 36''2$&$+35''3$&$+21''1$\\
März&6&$185\ 49\quad3.8$&$+5.0$&$+4.6$&$16\quad4\ 24.0$&$+34.5$&$+23.0$
\end{tabular}}
\end{center}

Dieser Vergleich zeigt nun deutlich, dass die Beobachtungen des Planeten gegen Ende März immer stärker von der Ellipse von \textsc{Gauß} abweichen werden. Selbst wenn sie im Juni um 5 Minuten abweichen sollten, was doch unwahrscheinlich ist, wäre das \textit{völlig belanglos}. Der Fehler kann sich nämlich nur allmählich anhäufen. Wenn man also jeweils die vorausgegangenen Beobachtungen mit den Elementen oder der Ephemeride vergleicht, wird man ihn stets mit. hinreichender Genauigkeit im Voraus kennen. Dr.\ \textsc{Gauß} schreibt uns auch, er habe bereits Vorbereitungen getroffen, um die Elemente sofort wieder so ändern zu können, dass sie sich den neuesten Beobachtungen aufs Genaueste anschließen, sobald sich der Fehler auf eine Minute anhäufen sollte. Dann kann man wieder für geraume Zeit eine sehr genaue Übereinstimmung erwarten. Sobald aber die Beobachtungen für dieses Jahr abgeschlossen sind und vollständig vorliegen, will er die wahre Ellipse mit allen ihren Störungsgleichungen so genau wie möglich bestimmen.
\par\nobreak\noindent\hfill[\textsc{von Zach}.]\par
